\section{函数的应用（二）}
\subsection{函数的增长速度}
\clearpage
\subsection{函数模型}
\clearpage

\hw
% 第一题：值域问题




% 第四题：生物丰富度指数选择题（短选项, 单行排版）
\begin{homework}
生物丰富度指数 \( d = \frac{S-1}{\ln N} \) 是河流水质的一个评价指标, 其中 \( S \), \( N \) 分别表示河流中的生物种类数与生物个体总数. 生物丰富度指数 \( d \) 越大, 水质越好. 如果某河流治理前后的生物种类数 \( S \) 没有变化, 生物个体总数由 \( N_{1} \) 变为 \( N_{2} \), 生物丰富度指数由 2.1 提高到 3.15, 则

\begin{table}[H]
  \centering
\begin{tabular*}{\hsize}{@{}@{\extracolsep{\fill}}llll@{}}
  \text{A.}  \( 3N_{2} = 2N_{1} \) & \text{B.}  \( 2N_{2} = 3N_{1} \) & \text{C.}  \( N_{2}^{2} = N_{1}^{3} \) & \text{D.}  \( N_{2}^{3} = N_{1}^{2} \)
\end{tabular*}
\end{table}
\end{homework}

% 第五题：天文学选择题（短选项, 单行排版）
\begin{homework}
在天文学中, 天体的明暗程度可以用星等或亮度来描述. 两颗星的星等与亮度满足
 \( m_{2} - m_{1} = \dfrac{5}{2} \lg \dfrac{E_{1}}{E_{2}} \), 其
 中星等为 \( m_{k} \) 的星的亮度为 \( E_{k} \)（k=1,2）. 已知太阳的星等是$ -26.7$
 , 天狼星的星等是 $-1.45$, 则太阳与天狼星的亮度的比值为

\begin{table}[H]
  \centering
\begin{tabular*}{\hsize}{@{}@{\extracolsep{\fill}}llll@{}}
  \text{A.} \( 10^{10.1} \) & \text{B.} \( 10.1 \) & \text{C.} \( \lg 10.1 \) & \text{D.} \( 10^{-10.1} \) % Note: User input has duplicated option for A and D, kept as provided
\end{tabular*}
\end{table}
\end{homework}

% 第六题：函数性质选择题（长选项, 分两行排版）


% =============== 第一张图片题目 ===============
\begin{homework}
假设某飞行器在空中高速飞行时所受的阻力 $f$ 满足公式 $ f = \frac{1}{2}\rho C S v^{2} $, 其中 $\rho$ 是空气密度, $S$ 是该飞行器的迎风面积, $v$ 是该飞行器相对于空气的速度, $C$ 是空气阻力系数, 功率 $P = fv$. 当 $\rho$, $S$ 不变, $v$ 比原来提高$10\%$时, 下列说法正确的是（ ）

% 短选项单行排版
\begin{table}[H]
  \centering
  \begin{tabular*}{\hsize}{@{}@{\extracolsep{\fill}}ll@{}}
    \text{A.} 若 $C$ 不变, 则 $P$ 比原来提高不超过 $30\%$  &\text{B.} 若 $C$ 不变, 则 $P$ 比原来提高超过 $40\%$  \\
    \text{C.} 为使 $P$ 不变, 则 $C$ 比原来降低不超过 $30\%$   &\text{D.} 为使 $P$ 不变, 则 $C$ 比原来降低超过 $40\%$\\
  \end{tabular*}
\end{table}
\end{homework}

\begin{homework}
德国心理学家研究发现, 人类大脑对事物的遗忘规律可用函数 $ y = 1 - 0.6 t^{0.27} $ 描述记忆率$y$随时间$t$（小时）的变化. 则记忆率为$50\%$时经过的时间约为（ ）

（参考数据：$\lg2 \approx 0.30$, $\lg3 \approx 0.48$）

\begin{table}[H]
  \centering
  \begin{tabular*}{\hsize}{@{}@{\extracolsep{\fill}}llll@{}}
    \text{A.} $2$ 小时 & 
    \text{B.} $0.8$ 小时 & 
    \text{C.} $0.5$ 小时 & 
    \text{D.} $0.2$ 小时
  \end{tabular*}
\end{table}
\end{homework}

\begin{homework}
某放射性物质的质量每年比前一年衰减$5\%$, 初始质量为$m_0$, $10$年后的质量为$m'$. 则与$\dfrac{m'}{m_{0}}$最接近的是（ ）

\begin{table}[H]
  \centering
  \begin{tabular*}{\hsize}{@{}@{\extracolsep{\fill}}llll@{}}
    \text{A.} $70\%$ & 
    \text{B.} $65\%$ & 
    \text{C.} $60\%$ & 
    \text{D.} $55\%$
  \end{tabular*}
\end{table}
\end{homework}

% =============== 第二张图片题目 ===============
\begin{homework}
某纯净水制造厂在净化水过程中, 每增加一次过滤可使水中杂质减少$50\%$. 要使水中杂质减少到原来的$10\%$以下, 则至少需要过滤（ ）

（参考数据：$\lg2 \approx 0.3010$）

\begin{table}[H]
  \centering
  \begin{tabular*}{\hsize}{@{}@{\extracolsep{\fill}}llll@{}}
    \text{A.} $2$ 次 & 
    \text{B.} $3$ 次 & 
    \text{C.} $4$ 次 & 
    \text{D.} $5$ 次
  \end{tabular*}
\end{table}
\end{homework}

\begin{homework}
垃圾分类投放可带来经济效益. 某市发放积分卡, 每分类投放1kg积1分; 若家庭月投放总量不低于100kg, 则额外奖励$x$分（$x \in \mathbb{N}^*$）. 积分按0.1元/分兑换. 
\begin{enumerate}
    \item 当$x=10$时, 若某家庭产生120kg垃圾, 兑换金额为\underline{\hspace{2cm}}元; 
    \item 为保证兑换金额不超过当月收益的$40\%$, 则$x$的最大值为\underline{\hspace{2cm}}. 
\end{enumerate}
\end{homework}

\begin{homework}
企业产量模型为$Q=AK^{\alpha}L^{\beta}$（$A>0$, $K>0$, $L>0$, $0<\alpha<1$, $0<\beta<1$）. 当$A$不变, $K$与$L$均变为原来的$2$倍时, 正确的是（ ）

% 长选项分行排版
\begin{table}[H]
\begin{tabular}{@{}@{\extracolsep{\fill}}llll@{}}
\text{A.}  存在 $\alpha<1/2$ 和 $\beta<1/2$, 使得 $Q$ 不变 \\ 
\text{B.}  存在 $\alpha>1/2$ 和 $\beta>1/2$, 使得 $Q$ 变为原来的 $2$ 倍 \\ 
\text{C.}  若 $\alpha=1/4$, 则 $Q$ 最多变为原来的 $2$ 倍 \\ 
\text{D.}  若 $\alpha^{2}+\beta^{2}=1/2$, 则 $Q$ 最多变为原来的 $2$ 倍 \\ 
\end{tabular}
\end{table}
\end{homework}

% =============== 第三张图片题目 ===============
\begin{homework}
某工厂废气污染物含量$P$（单位：mg/L）与时间$t$（单位：h）满足$P=P_0 e^{-kt}$（$P_0,k>0$）. 若前$10$h污染物减少$19\%$, 则再过$5$h后污染物剩余（ ）

\begin{table}[H]
  \centering
  \begin{tabular*}{\hsize}{@{}@{\extracolsep{\fill}}llll@{}}
    \text{A.} $40.5\%$ & 
    \text{B.} $54\%$ & 
    \text{C.} $65.6\%$ & 
    \text{D.} $72.9\%$
  \end{tabular*}
\end{table}
\end{homework}

\begin{homework}某公司通过劳动统计分析发现, 
工人工作效率\(E\)与工作年限\(r\)、劳累程度\(T\)、劳动动机\(b\)相关,
并建立数学模型$E=10-10T\cdot b^{-0.14r}$（$r>0$, $0<T<1$, $1<b<5$）. 
对员工甲乙的下列结论：
\begin{enumerate}
  \item 甲乙劳动动机相同, 甲工作年限长且劳累程度弱, 则甲效率高
  \item 甲乙劳累程度相同, 甲工作年限长且劳动动机高, 则甲效率高
  \item 甲乙工作年限相同, 甲效率高且劳动动机低, 则甲劳累程度强
  \item 甲乙劳动动机相同, 甲效率高且工作年限短, 则甲劳累程度弱
\end{enumerate}
其中所有正确结论的序号是\underline{\hspace{2cm}}. 
\end{homework}

\begin{homework}
火箭最大速度$v$（单位：km/s）满足$v=2000\ln(1+\frac{M}{m})$. 当燃料比$t_0=\frac{M}{m}$时速度$v=v_0$. 若使$v=2v_0$, 则燃料比应为（ ）

\begin{table}[H]
  \centering
  \begin{tabular*}{\hsize}{@{}@{\extracolsep{\fill}}llll@{}}
    \text{A.} $2t_{0}^{2}$ & 
    \text{B.} $t_{0}^{2}+t_{0}$ & 
    \text{C.} $2t_{0}$ & 
    \text{D.} $t_{0}^{2}+2t_{0}$
  \end{tabular*}
\end{table}
\end{homework}

\begin{homework}
火箭最大速度满足$v=2\ln(1+\frac{M}{N})$. 若$v=12\mathrm{km/s}$, 则$\dfrac{M}{N}$最接近的是（ ）

（参考数据：$\e\approx 2.71828$）

\begin{table}[H]
  \centering
  \begin{tabular*}{\hsize}{@{}@{\extracolsep{\fill}}llll@{}}
    \text{A.} $200$ & 
    \text{B.} $400$ & 
    \text{C.} $600$ & 
    \text{D.} $800$
  \end{tabular*}
\end{table}
\end{homework}